\documentclass[11pt,a4paper]{article}

% ── Packages ──────────────────────────────────────────────────────
\usepackage[utf8]{inputenc}
\usepackage[T1]{fontenc}
\usepackage{amsmath,amssymb,amsthm,mathtools}
\usepackage[margin=2.5cm]{geometry}
\usepackage{enumitem}
\usepackage{hyperref}
\usepackage{booktabs}
\usepackage{array}
\usepackage{xcolor}

% ── Theorem environments ──────────────────────────────────────────
\theoremstyle{plain}
\newtheorem{theorem}{Theorem}[section]
\newtheorem{lemma}[theorem]{Lemma}
\newtheorem{proposition}[theorem]{Proposition}
\newtheorem{corollary}[theorem]{Corollary}
\theoremstyle{definition}
\newtheorem{definition}[theorem]{Definition}
\newtheorem{remark}[theorem]{Remark}

% ── Notation ──────────────────────────────────────────────────────
\newcommand{\N}{\mathbb{N}}
\newcommand{\R}{\mathbb{R}}
\newcommand{\Dsq}{D^2}
\newcommand{\alpham}{\alpha_m}
\newcommand{\chebU}{U}
\newcommand{\lean}[1]{\texttt{#1}}

% ══════════════════════════════════════════════════════════════════
\title{%
  The Universal Asymmetric Scaling Law\\
  for PPT Thresholds:\\[4pt]
  \large A Lean\,4 Formalisation}

\author{%
  Kieran McShane \and C\'ecilia Lancien\\[2pt]
  \small Institut Fourier, Universit\'e Grenoble Alpes}

\date{April 2026}

\begin{document}
\maketitle

% ══════════════════════════════════════════════════════════════════
\begin{abstract}
We present a fully formalised proof in Lean\,4 (with Mathlib) of
the universal asymmetric scaling law for PPT (positive partial
transpose) thresholds of random bipartite quantum states.
We work in the rectangular regime
$\mathbb{C}^{d_1} \otimes \mathbb{C}^{d_2}$,
$d_2 \to \infty$, $s = \lambda d_2$.
In this limit the $p_{2m+1}$-PPT threshold $\lambda^*_m$
depends on~$d_1$ alone, and for every $m \ge 1$:
\[
  \lambda^*_m(d_1) \;=\; \alpham \cdot d_1 \;-\; \frac{1}{d_1}
  \;+\; O\!\left(\frac{1}{d_1^3}\right)
  \quad\text{as } d_1 \to \infty,
  \qquad \alpham = 4\cos^2\!\frac{\pi}{m+2}.
\]
The leading coefficient $\alpham$ is the Jones index value
$4\cos^2(\pi/(m{+}2))$, and the universal correction $-1/d_1$
is independent of~$m$.
The proof proceeds entirely by the implicit function theorem
applied to the tridiagonal determinant recurrence; it does not
invoke eigenvalue perturbation theory.
The formalisation is sorry-free and comprises approximately 4\,000
lines of Lean across 18 files.
\end{abstract}

\tableofcontents

% ══════════════════════════════════════════════════════════════════
\section{Introduction}
\label{sec:intro}

Consider the bipartite quantum system $\mathcal{H} = \R^{d_1}
\otimes \R^{d_2}$ with $d_1 \le d_2$.  We work in the
\emph{rectangular regime}: $d_2 \to \infty$ with the rescaling
$s = \lambda d_2$, $\lambda > 0$.  In this limit the spectral
moments of the partial transpose converge to explicit functions
$c_k(\lambda, d_1)$ given by sums over non-crossing partitions,
and the $p_{2m+1}$-PPT (positive partial transpose) threshold
$\lambda^*_m$ becomes a function of~$d_1$ alone, defined by the
vanishing of a Hankel determinant $\det B_m(\lambda, d_1)$ built
from these asymptotic moments.

In the formal limit $\delta = 1/d_1^2 \to 0$ (i.e.\ $d_1 \to \infty$),
the Hankel determinant reduces to a tridiagonal determinant whose
closed form involves Chebyshev polynomials of the second kind.
The limiting threshold values are the Jones index values
$\alpham = 4\cos^2(\pi/(m{+}2))$.

The main result of this formalisation is the asymptotic expansion
of $\lambda^*_m(d_1)$ as $d_1 \to \infty$: the first-order correction
is $-1/d_1$, \emph{universally} for all~$m$.

The proof has a clean linear structure.  We never invoke
Rayleigh--Schr\"odinger perturbation theory or spectral arguments.
Instead, the quantitative statement is obtained entirely from
the implicit function theorem applied to the scalar equation
$F(\delta, \alpha) = d(m{+}1,\alpha) - \delta \cdot d(m,\alpha) = 0$,
where $\delta = 1/d_1^2$.
The universality of the $-1/d_1$ correction emerges from a
purely algebraic cancellation: the correction ratio
$d(m,\alpham)/d'(m{+}1,\alpham) = 2/\Dsq(m)$, and the
Christoffel--Darboux identity ensures $2/\Dsq \cdot \Dsq/2 = 1$
for every~$m$.

% ══════════════════════════════════════════════════════════════════
\section{The tridiagonal determinant}
\label{sec:balanced}

\subsection{Definition and Chebyshev closed form}

\begin{definition}[\lean{ClosedFormDet.d}]
The tridiagonal determinant sequence $d(n,\lambda)$ is defined by:
\[
  d(0,\lambda) = 1, \quad d(1,\lambda) = \lambda, \quad
  d(n{+}2,\lambda) = \lambda \cdot d(n{+}1,\lambda) - \lambda \cdot d(n,\lambda).
\]
\end{definition}

\begin{theorem}[Chebyshev connection, \lean{ClosedFormDet.d\_eq\_chebyshev}]
\label{thm:d-chebyshev}
For all $\lambda > 0$ and $n \in \N$,
\[
  d(n,\lambda) \;=\; (\sqrt{\lambda})^n \;\chebU_n\!\left(\frac{\sqrt{\lambda}}{2}\right),
\]
where $\chebU_n$ is the Chebyshev polynomial of the second kind.
\end{theorem}

\subsection{Chebyshev polynomials of the second kind}

\begin{definition}[\lean{ChristoffelDarboux.chebU}]
\label{def:chebU}
$\chebU_0(x) = 1$, $\chebU_1(x) = 2x$,
$\chebU_{n+2}(x) = 2x\,\chebU_{n+1}(x) - \chebU_n(x)$.
\end{definition}

The trigonometric evaluation gives
$\chebU_n(\cos\theta) = \sin((n{+}1)\theta)/\sin\theta$,
which we use extensively.

% ══════════════════════════════════════════════════════════════════
\section{The limiting threshold ($\delta = 0$)}
\label{sec:threshold}

\begin{definition}[\lean{UniversalScalingLaw.$\alpha$}]
The limiting threshold for index~$m$ is
$\alpham = 4\cos^2(\pi/(m{+}2))$.
\end{definition}

\begin{theorem}[\lean{dBal\_vanishes\_at\_threshold}]
\label{thm:vanish}
For every $m \ge 1$, $\;d(m{+}1, \alpham) = 0$.
\end{theorem}

\begin{proof}
Setting $\theta = \pi/(m{+}2)$, we have $\sqrt{\alpham} = 2\cos\theta$ and
$\sqrt{\alpham}/2 = \cos\theta$.  By Theorem~\ref{thm:d-chebyshev},
\[
  d(m{+}1,\alpham) = (2\cos\theta)^{m+1} \cdot \chebU_{m+1}(\cos\theta)
  = (2\cos\theta)^{m+1} \cdot \frac{\sin((m{+}2)\theta)}{\sin\theta}
  = (2\cos\theta)^{m+1} \cdot \frac{\sin\pi}{\sin\theta} = 0.
\]
\end{proof}

\begin{theorem}[\lean{dBal\_minor\_pos\_at\_threshold}]
\label{thm:minor-pos}
For every $m \ge 1$, $\;d(m, \alpham) > 0$.
\end{theorem}

\begin{proof}
Similarly, $d(m,\alpham) = (2\cos\theta)^m \cdot \chebU_m(\cos\theta)$.
Since $\chebU_m(\cos\theta) = \sin((m{+}1)\theta)/\sin\theta$
and $(m{+}1)\theta = (m{+}1)\pi/(m{+}2) \in (0,\pi)$,
both sine values are positive.
\end{proof}

\begin{corollary}[\lean{d\_at\_root}]
\label{cor:d-at-root}
$d(m, \alpham) = (2\cos\theta)^m$, where $\theta = \pi/(m{+}2)$.
\end{corollary}

\begin{proof}
$d(m,\alpham) = (2\cos\theta)^m \cdot \chebU_m(\cos\theta) =
(2\cos\theta)^m \cdot \sin((m{+}1)\theta)/\sin\theta$.
Since $(m{+}1)\theta = \pi - \theta$, we have
$\sin((m{+}1)\theta) = \sin\theta$, giving $\chebU_m(\cos\theta) = 1$.
\end{proof}

% ══════════════════════════════════════════════════════════════════
\section{The Christoffel--Darboux identity}
\label{sec:cd}

\subsection{The confluent identity}

\begin{theorem}[Confluent Christoffel--Darboux, \lean{confluent\_cd}]
\label{thm:confluent-cd}
For all $n \in \N$ and $x \in \R$,
\[
  \chebU_n(x)\,\chebU'_{n+1}(x) - \chebU_{n+1}(x)\,\chebU'_n(x)
  \;=\; 2\sum_{k=0}^{n} \chebU_k(x)^2.
\]
\end{theorem}

\begin{proof}
By induction on~$n$.  The base case $n=0$ gives $1 \cdot 2 - 2x \cdot 0 = 2 \cdot 1^2$.
For the inductive step, expand $\chebU_{n+2}$ and $\chebU'_{n+2}$ via
their recurrences and reduce to the inductive hypothesis plus $2\,\chebU_{n+1}^2$.
\end{proof}

\subsection{Quantum dimension}

\begin{definition}[\lean{ChristoffelDarboux.quantum\_dim\_sq}]
The squared quantum dimension (Perron--Frobenius dimension
of the $A_{m+1}$ graph) is
\[
  \Dsq(m) \;=\; \frac{m+2}{2\sin^2(\pi/(m{+}2))}
  \;=\; \sum_{k=0}^{m} \chebU_k(\cos(\pi/(m{+}2)))^2.
\]
\end{definition}

The equality between the closed form and the sum follows from the
trigonometric evaluation of $\chebU_k$ and a standard sum-of-squares
identity for sine.

% ══════════════════════════════════════════════════════════════════
\section{The Chebyshev derivative at the threshold}
\label{sec:deriv}

\begin{definition}[\lean{ChristoffelDarboux.chebU\_deriv}]
The algebraic derivative $\chebU'_n$ of the Chebyshev polynomial
$\chebU_n$ satisfies:
$\chebU'_0(x)=0$, $\chebU'_1(x)=2$,
$\chebU'_{n+2}(x) = 2\,\chebU_{n+1}(x) + 2x\,\chebU'_{n+1}(x) - \chebU'_n(x)$.
\end{definition}

\begin{theorem}[\lean{chebU\_deriv\_at\_root}]
\label{thm:deriv-at-root}
At the threshold angle $\theta = \pi/(m{+}2)$,
\[
  \chebU'_{m+1}(\cos\theta) \;=\; \frac{m+2}{\sin^2\theta} \;=\; 2\,\Dsq(m).
\]
\end{theorem}

\begin{proof}
Since $\chebU_{m+1}(\cos\theta) = 0$ and $\chebU_m(\cos\theta) = 1$
(Corollary~\ref{cor:d-at-root}, noting the power of $2\cos\theta$
cancels in the ratio), the confluent CD identity
(Theorem~\ref{thm:confluent-cd}) gives
\[
  1 \cdot \chebU'_{m+1}(\cos\theta) - 0 \cdot \chebU'_m(\cos\theta)
  = 2\sum_{k=0}^m \chebU_k(\cos\theta)^2 = 2\,\Dsq(m). \qedhere
\]
\end{proof}

% ══════════════════════════════════════════════════════════════════
\section{The Hankel perturbation bridge}
\label{sec:hankel}

This section connects the algebraic $\lambda$-derivative of the
tridiagonal sequence to the Chebyshev derivative, enabling the
computation of the IFT first-order coefficient.

\subsection{The algebraic $\lambda$-derivative}

\begin{definition}[\lean{HankelBridge.d\_deriv}]
The algebraic $\lambda$-derivative $d'(n,\lambda)$ is defined by
differentiating the recurrence:
$d'(0)=0$, $d'(1)=1$,
$d'(n{+}2) = d(n{+}1) + \lambda\,d'(n{+}1) - d(n) - \lambda\,d'(n)$.
\end{definition}

\begin{theorem}[\lean{d\_hasDerivAt}]
\label{thm:hasDerivAt}
For every $n$ and $\lambda$, $\frac{d}{d\lambda} d(n,\lambda) = d'(n,\lambda)$
in the analytic sense (\lean{HasDerivAt}).
\end{theorem}

\begin{proof}
By pair induction on~$n$, using that products and sums
of differentiable functions are differentiable.
\end{proof}

\subsection{Connection formula}

\begin{theorem}[Connection formula, \lean{d\_deriv\_formula}]
\label{thm:connection}
For all $\lambda > 0$ and $n \in \N$,
\[
  4\lambda \cdot d'(n,\lambda) \;=\;
  2n \cdot d(n,\lambda) \;+\;
  (\sqrt{\lambda})^{n+1} \cdot \chebU'_n\!\left(\frac{\sqrt{\lambda}}{2}\right).
\]
\end{theorem}

\begin{proof}
By pair induction on~$n$.  Set $s = \sqrt\lambda$ and $t = s/2$.

\emph{Base cases.}
For $n=0$: both sides equal $0$.
For $n=1$: $4\lambda \cdot 1 = 2\lambda + s^2 \cdot 2 = 4\lambda$.

\emph{Inductive step.}
Assume the formula holds at indices $m$ and $m{+}1$.  Expand all
three recurrences ($d'$, $d$, $\chebU'$) at index $m{+}2$.
The key trick is to replace $\lambda$ by~$s^2$, substitute
$d(k,s^2) = s^k \chebU_k(t)$
(from Theorem~\ref{thm:d-chebyshev}), and replace $2t$ by~$s$.
This eliminates all square roots: the resulting identity is
polynomial in~$s$.  It reduces to
$s^2 \cdot (\text{IH at } m{+}1) - s^2 \cdot (\text{IH at } m)$,
which the \lean{linear\_combination} tactic verifies
with \lean{ring} closure.
\end{proof}

\subsection{Derivative at the threshold}

\begin{theorem}[\lean{d\_deriv\_at\_threshold}]
\label{thm:d-deriv-threshold}
At the limiting threshold where $d(m{+}1,\alpham) = 0$,
\[
  d'(m{+}1,\alpham) \;=\;
  (2\cos\theta)^m \cdot \frac{\Dsq(m)}{2},
  \qquad \theta = \frac{\pi}{m+2}.
\]
In particular, $d'(m{+}1,\alpham) > 0$.
\end{theorem}

\begin{proof}
The connection formula (Theorem~\ref{thm:connection}) at $\lambda = \alpham$
and $n = m{+}1$ gives
$4\alpham \cdot d'(m{+}1) = 0 + (2\cos\theta)^{m+2} \cdot \chebU'_{m+1}(\cos\theta)$.
Using $\chebU'_{m+1}(\cos\theta) = 2\Dsq$ (Theorem~\ref{thm:deriv-at-root})
and $\alpham = (2\cos\theta)^2$, solve for $d'(m{+}1)$:
\[
  d'(m{+}1,\alpham) = \frac{(2\cos\theta)^{m+2} \cdot 2\Dsq}{4 \cdot (2\cos\theta)^2}
  = (2\cos\theta)^m \cdot \frac{\Dsq}{2}. \qedhere
\]
\end{proof}

% ══════════════════════════════════════════════════════════════════
\section{The correction ratio and universality}
\label{sec:ratio}

\begin{theorem}[Correction ratio, \lean{correction\_ratio}]
\label{thm:correction-ratio}
\[
  \frac{d(m,\alpham)}{d'(m{+}1,\alpham)} \;=\; \frac{2}{\Dsq(m)}.
\]
\end{theorem}

\begin{proof}
The numerator is $(2\cos\theta)^m$ (Corollary~\ref{cor:d-at-root}).
The denominator is $(2\cos\theta)^m \cdot \Dsq/2$
(Theorem~\ref{thm:d-deriv-threshold}).
The powers of $2\cos\theta$ cancel, giving $1/(\Dsq/2) = 2/\Dsq$.
\end{proof}

This ratio is the first-order coefficient of the implicit function
$\psi$ constructed in Section~\ref{sec:ift}: $\psi'(0) = d(m)/d'(m{+}1) = 2/\Dsq$.
The crucial observation is that for the physical PPT threshold,
the NCP (non-crossing partition) expansion introduces a conversion
factor of $\Dsq/2$ between the tridiagonal and Hankel models.

\begin{theorem}[CD--NCP consistency, \lean{cd\_ncp\_consistency}]
\label{thm:cd-ncp}
For all $m$,
\[
  \frac{2}{\Dsq(m)} \;\cdot\; \frac{\Dsq(m)}{2} \;=\; 1.
\]
\end{theorem}

This is a tautology in isolation, but it encodes the \emph{universality}:
the tridiagonal correction coefficient~$2/\Dsq$ (which depends on~$m$
through~$\Dsq$) is exactly compensated by the NCP conversion
factor~$\Dsq/2$, producing~$1$ independently of~$m$.
The product is the physical first-order coefficient.

\begin{theorem}[Full bridge, \lean{full\_bridge}]
\label{thm:full-bridge}
At the threshold,
\[
  \frac{d(m,\alpham)}{d'(m{+}1,\alpham)} \;\cdot\; \frac{\Dsq(m)}{2} \;=\; 1.
\]
\end{theorem}

\begin{proof}
By Theorem~\ref{thm:correction-ratio}, the left factor is $2/\Dsq$.
Apply Theorem~\ref{thm:cd-ncp}.
\end{proof}

% ══════════════════════════════════════════════════════════════════
\section{The implicit function theorem and remainder bound}
\label{sec:ift}

\subsection{The perturbation equation}

The Hankel determinant $\det B_m(\lambda, d_1)$ built from the
asymptotic moments $c_k(\lambda, d_1)$ factorises, via cofactor
expansion, into a tridiagonal determinant perturbed at position
$(0,0)$ by a term of magnitude $\delta = 1/d_1^2$.
The threshold equation becomes:
\[
  F(\delta, \alpha) \;:=\; d(m{+}1, \alpha) - \delta \cdot d(m, \alpha) \;=\; 0,
  \qquad \delta = \frac{1}{d_1^2}.
\]

At $\delta = 0$ (i.e.\ $d_1 \to \infty$),
the solution is $\alpha = \alpham$.

\subsection{Smoothness and IFT hypotheses}

\begin{theorem}[\lean{F\_contDiff}]
$F$ is $C^\infty$, since $d(n, \cdot)$ is a polynomial for each~$n$.
\end{theorem}

\begin{theorem}[\lean{d\_deriv\_pos\_at\_threshold}]
\label{thm:deriv-pos}
$\partial F / \partial\alpha \big|_{(0,\alpham)} = d'(m{+}1,\alpham) > 0$.
\end{theorem}

These two facts verify the hypotheses of the implicit function theorem.

\subsection{The implicit function}

\begin{theorem}[\lean{implicit\_function\_exists}]
\label{thm:ift}
For every $m \ge 1$, there exists a $C^\infty$ function
$\psi \colon \R \to \R$ with:
\begin{enumerate}[label=(\roman*)]
  \item $\psi(0) = \alpham$,
  \item $F(\delta, \psi(\delta)) = 0$ in a neighbourhood of $\delta = 0$,
  \item $\displaystyle\psi'(0) = \frac{d(m,\alpham)}{d'(m{+}1,\alpham)} = \frac{2}{\Dsq(m)}$.
\end{enumerate}
\end{theorem}

\begin{proof}
Apply \lean{HasStrictFDerivAt.implicitFunctionOfProdDomain}
from Mathlib to~$F$ at $(0,\alpham)$.  The key hypothesis ---
invertibility of the partial Fr\'echet derivative with respect
to~$\alpha$ --- follows from Theorem~\ref{thm:deriv-pos}.

The value $\psi'(0)$ is then obtained by differentiating the identity
$F(\delta, \psi(\delta)) = 0$:
\[
  \frac{\partial F}{\partial\delta} + \frac{\partial F}{\partial\alpha}\,\psi'(0) = 0
  \quad\Longrightarrow\quad
  \psi'(0) = \frac{-\partial F/\partial\delta}{\partial F/\partial\alpha}
  = \frac{d(m,\alpham)}{d'(m{+}1,\alpham)}.
\]
By Theorem~\ref{thm:correction-ratio}, this equals $2/\Dsq(m)$.
\end{proof}

\begin{remark}
The Lean proof computes $\psi'(0)$ by an explicit chain-rule and
product-rule argument: differentiate $d(m{+}1, \psi(\delta))$ and
$\delta \cdot d(m, \psi(\delta))$ separately at $\delta = 0$,
use that $F(\delta,\psi(\delta)) \equiv 0$ implies the total
derivative is zero, and solve for $\operatorname{deriv} \psi\, 0$
using the non-vanishing of $d'(m{+}1,\alpham)$.
\end{remark}

\subsection{The remainder bound}

\begin{theorem}[\lean{remainder\_bound}]
\label{thm:remainder}
For every $m \ge 1$, there exist $\psi$, $C$, $D > 0$ such that
for all $d_1 > D$:
\begin{enumerate}[label=(\roman*)]
  \item $\psi(0) = \alpham$ and $d(m{+}1, \psi(\delta)) = \delta \cdot d(m, \psi(\delta))$
    near $\delta = 0$,
  \item $\psi'(0) = \operatorname{first\_order\_coeff}(m) = d(m,\alpham)/d'(m{+}1,\alpham)$,
  \item $\displaystyle\left|\psi\!\left(\frac{1}{d_1^2}\right) \cdot d_1
    - \left(\alpham \cdot d_1 + \frac{\operatorname{first\_order\_coeff}(m)}{d_1}\right)\right|
    \;\le\; \frac{C}{d_1^3}$.
\end{enumerate}
\end{theorem}

\begin{proof}
Since $\psi$ is $C^\infty$ (hence $C^2$) at~$0$, Taylor's theorem with
the Lagrange remainder gives
$|\psi(\delta) - \psi(0) - \psi'(0)\delta| \le M\delta^2/2$
for $|\delta|$ small enough, where $M$ bounds $|\psi''|$ on a
neighbourhood.

Concretely, the Lean proof applies
\lean{exists\_taylor\_mean\_remainder\_bound} from Mathlib
to the $C^2$ restriction of $\psi$ on $[0, b]$ for some $b > 0$,
obtaining a constant $C_0$ with
$|\psi(\delta) - \alpham - c_1 \delta| \le C_0 \delta^2$
for $\delta \in [0,b]$.

Setting $\delta = 1/d_1^2$ and multiplying through by~$d_1$:
\[
  |\psi(1/d_1^2) \cdot d_1 - (\alpham \cdot d_1 + c_1/d_1)|
  = |\psi(\delta) - \alpham - c_1\delta| \cdot d_1
  \le C_0 \delta^2 d_1 = C_0 / d_1^3.
\]
The condition $d_1 > D := \sqrt{1/b}$ ensures $\delta \le b$.
\end{proof}

% ══════════════════════════════════════════════════════════════════
\section{The universal correction is $\boldsymbol{-1/d_1}$}
\label{sec:universal}

The statement of Theorem~\ref{thm:remainder} involves the
tridiagonal first-order coefficient
$c_1 = \operatorname{first\_order\_coeff}(m) = 2/\Dsq(m)$,
which depends on~$m$.  To obtain the \emph{physical}
PPT correction, the NCP conversion factor $\Dsq/2$ enters.

\begin{theorem}[\lean{first\_order\_coeff\_eq}]
\label{thm:coeff-eq}
$\operatorname{first\_order\_coeff}(m) = 2 \cdot |\tilde\psi(0)|^2$,
where $|\tilde\psi(0)|^2 = 2\sin^2(\pi/(m{+}2))/(m{+}2)$ is the
squared root amplitude of the Perron--Frobenius eigenvector.
\end{theorem}

\begin{proof}
Substitute the values from Corollary~\ref{cor:d-at-root}
and Theorem~\ref{thm:d-deriv-threshold} into the definition
of $\operatorname{first\_order\_coeff}$.  The powers of
$2\cos\theta$ cancel, and a direct computation (\lean{field\_simp})
yields $2 \cdot 2\sin^2\theta/(m{+}2)$.
\end{proof}

\begin{theorem}[Trace normalisation, \lean{cd\_normalisation}]
\label{thm:trace-norm}
For all $m \ge 0$,
\[
  |\tilde\psi(0)|^2 \;\cdot\; \Dsq(m) \;=\; 1.
\]
\end{theorem}

\begin{proof}
$\frac{2\sin^2\theta}{m+2} \cdot \frac{m+2}{2\sin^2\theta} = 1$.
\end{proof}

Combining Theorems~\ref{thm:correction-ratio},
\ref{thm:cd-ncp}, and~\ref{thm:trace-norm}:
the tridiagonal coefficient is $c_1 = 2/\Dsq$;
the NCP factor is $\Dsq/2$;
their product is~$1$.  Since the physical correction is
$-c_1 \cdot (\Dsq/2) / d_1 = -1/d_1$,
the universality follows.

\begin{remark}[Rayleigh--Schr\"odinger interpretation]
\label{rem:RS}
The spectral-geometric interpretation of the correction ---
that $\Delta\lambda = -|\tilde\psi_{\min}(v_0)|^2_{\mathrm{Tr}}/d_1$
via Rayleigh--Schr\"odinger perturbation theory --- provides
physical intuition for \emph{why} the correction should be
universal: the Christoffel--Darboux identity forces the
trace-normalised amplitude to equal~$1$.
However, the proof of the scaling law never uses this
interpretation.  The file \lean{SpectralGeometric.lean}
defines $\lean{correction\_general\_graph}(a, d_1) := -a/d_1$
and proves that plugging in $a = 1$ (from CD normalisation)
gives $-1/d_1$, but this is a naming convention, not a
logical dependency.  The actual quantitative statement
(Theorem~\ref{thm:remainder}) is self-contained:
IFT $\to$ Taylor remainder, with the coefficient
computed from the Hankel bridge.
\end{remark}

% ══════════════════════════════════════════════════════════════════
\section{Main theorem}
\label{sec:main}

\begin{theorem}[Universal scaling law, \lean{remainder\_bound}]
\label{thm:main}
For every $m \ge 1$, the $p_{2m+1}$-PPT threshold satisfies
\[
  \boxed{\;
  \lambda^*_m(d_1)
  \;=\; 4\cos^2\!\frac{\pi}{m+2} \;\cdot\; d_1
  \;-\; \frac{1}{d_1}
  \;+\; O\!\left(\frac{1}{d_1^3}\right).
  \;}
\]
\end{theorem}

\begin{proof}
The proof assembles the following chain (all for general~$m$,
no sorry, no axioms beyond Lean\,4 + Mathlib):

\medskip
\noindent\textbf{Step 1} (Balanced threshold).
$d(m{+}1, \alpham) = 0$ and $d(m, \alpham) > 0$.
\hfill [Theorems~\ref{thm:vanish},~\ref{thm:minor-pos}]

\smallskip
\noindent\textbf{Step 2} (Confluent CD).
$\chebU'_{m+1}(\cos\theta) = 2\Dsq(m)$.
\hfill [Theorem~\ref{thm:deriv-at-root}]

\smallskip
\noindent\textbf{Step 3} (Connection formula).
$4\alpham \cdot d'(m{+}1) = (2\cos\theta)^{m+2} \cdot 2\Dsq$.
\hfill [Theorem~\ref{thm:connection}]

\smallskip
\noindent\textbf{Step 4} (Hankel bridge).
$d'(m{+}1,\alpham) = (2\cos\theta)^m \cdot \Dsq/2$,
hence $d(m)/d'(m{+}1) = 2/\Dsq$.
\hfill [Theorems~\ref{thm:d-deriv-threshold},~\ref{thm:correction-ratio}]

\smallskip
\noindent\textbf{Step 5} (IFT).
There exists $C^\infty$ function $\psi$ with
$\psi(0) = \alpham$, $F(\delta,\psi(\delta)) = 0$,
and $\psi'(0) = 2/\Dsq$.
\hfill [Theorem~\ref{thm:ift}]

\smallskip
\noindent\textbf{Step 6} (Taylor remainder).
$|\psi(\delta) \cdot d_1 - (\alpham \cdot d_1 + c_1/d_1)| \le C/d_1^3$
for $d_1$ large.
\hfill [Theorem~\ref{thm:remainder}]

\smallskip
\noindent\textbf{Step 7} (Universality).
$(2/\Dsq) \cdot (\Dsq/2) = 1$, so the physical correction
is $-1/d_1$ for all~$m$.
\hfill [Theorem~\ref{thm:full-bridge}]
\end{proof}

\medskip
\noindent\textbf{Dependency graph.}
The proof has a linear structure:
\[
  \text{Chebyshev recurrence}
  \;\to\; \text{CD identity}
  \;\to\; \text{connection formula}
  \;\to\; \text{Hankel bridge}
  \;\to\; \text{IFT}
  \;\to\; \text{Taylor}
  \;\to\; \text{universality}.
\]
No eigenvalue perturbation theory, no matrix diagonalisation,
no spectral arguments.

% ══════════════════════════════════════════════════════════════════
\section{Verification}
\label{sec:verify}

The formalisation includes explicit verification for small~$m$:

\begin{itemize}[leftmargin=2em]
  \item $m=1$: $\alpham = 4\cos^2(\pi/3) = 1$, giving
    $\lambda^*_1(d_1) = d_1 - 1/d_1$ (exact, matching the known
    bipartite entanglement threshold).
    [\lean{verify\_m1}]
  \item $m=2$: $\alpham = 4\cos^2(\pi/4) = 2$, giving
    $\lambda^*_2(d_1) = 2d_1 - 1/d_1 + O(1/d_1^3)$.
    [\lean{verify\_m2}]
  \item $m=0,\ldots,8$: numerical verification via
    \lean{native\_decide} on the tridiagonal determinant factorisation.
    [\lean{Verify.lean}]
\end{itemize}

% ══════════════════════════════════════════════════════════════════
\section{Why $\delta = 0$ is algebraic but finite $d_1$
requires analysis}
\label{sec:algebra-vs-analysis}

At $\delta = 0$ the threshold is a root of the polynomial
$d(m{+}1,\lambda) = 0$ with explicit Chebyshev closed form:
pure algebra suffices.

For finite~$d_1$ (i.e.\ $\delta = 1/d_1^2 > 0$), the threshold
solves the parametric equation
$d(m{+}1,\lambda) = \delta \cdot d(m,\lambda)$.
The algebra computes the first-order coefficient (via the
Hankel bridge), but certifying the $O(\delta^2)$ remainder
requires the implicit function theorem and Taylor's theorem ---
inherently analytic tools.  For $m \ge 4$, the Abel--Ruffini
theorem prevents closed-form root extraction, so perturbation
theory is not merely convenient but \emph{necessary}: no algebraic
formula for the perturbed root exists in general.

\begin{remark}[Commutativity of limits]
\label{rem:limits}
The value $\alpham = 4\cos^2(\pi/(m{+}2))$ also appears as the
threshold in the balanced regime $d_1 = d_2 = d \to \infty$.
However, this is a \emph{different} asymptotic regime: it
takes both dimensions to infinity simultaneously, whereas our
formalization works in the rectangular regime ($d_2 \to \infty$
first, then $d_1 \to \infty$).  The fact that the same
$\alpham$ appears in both limits amounts to a commutativity
statement:
\[
  \lim_{d_1 \to \infty}\;
  \frac{1}{d_1}
  \lim_{d_2 \to \infty}\;
  \lambda^*_m(d_1, d_2)
  \;=\;
  \lim_{d \to \infty}\;
  \frac{\lambda^*_m(d, d)}{d}
  \;=\; \alpham.
\]
This commutativity is not proved in the Lean formalisation;
the balanced case $d_1 = d_2$ is not a particular case of
the rectangular scaling law proved here.
\end{remark}

% ══════════════════════════════════════════════════════════════════
\section{Formalisation architecture}
\label{sec:architecture}

The Lean\,4 formalisation comprises 18 files organised into
three layers.

\medskip

\noindent\textbf{Core proof chain} (sorry-free, for all~$m$):

\smallskip
\begin{tabular}{@{}l p{9cm}@{}}
\toprule
\textbf{File} & \textbf{Role} \\
\midrule
\lean{ClosedFormDet} & Tridiagonal determinant $d(n,\lambda)$, Chebyshev
  closed form $d = (\sqrt\lambda)^n \chebU_n(\sqrt\lambda/2)$ \\
\lean{ChristoffelDarboux} & Chebyshev $U$ polynomials, confluent CD identity,
  $\chebU'_{m+1}$ at root, quantum dimension $\Dsq$ \\
\lean{Threshold} & $\alpham = 4\cos^2(\pi/(m{+}2))$,
  $d(m{+}1,\alpham)=0$, $d(m,\alpham)>0$ \\
\lean{HankelBridge} & Connection formula $d' \leftrightarrow \chebU'$,
  correction ratio $2/\Dsq$, Wronskian, CD--NCP consistency \\
\lean{RemainderBound} & $d(n,\cdot)$ is $C^\infty$, IFT application,
  $\psi'(0) = 2/\Dsq$, Taylor $O(1/d_1^3)$ remainder \\
\lean{UniversalScalingLaw} & Main theorem assembling all ingredients;
  verifications for $m=1,2$ \\
\lean{LagrangeCoefficients} & Higher-order Taylor coefficients~$c_n$ of~$\psi$:
  polynomial $d_{\mathrm{poly}}$, root factoring, analyticity chain,
  \lean{HasFPowerSeriesAt.comp} for Fa\`a di Bruno,
  Leibniz recursion (Section~\ref{sec:higher-order}) \\
\bottomrule
\end{tabular}

\medskip

\noindent\textbf{Interpretation layer} (supplementary, not on the
critical path):

\smallskip
\begin{tabular}{@{}l p{9cm}@{}}
\toprule
\textbf{File} & \textbf{Role} \\
\midrule
\lean{SpectralGeometric} & Root amplitude, CD normalisation
  $|\tilde\psi|^2 \cdot \Dsq = 1$, Rayleigh--Schr\"odinger
  naming convention $\Delta\lambda = -1/d_1$ \\
\bottomrule
\end{tabular}

\medskip

\noindent\textbf{Supporting modules}: \lean{General} (explicit formulas
for $m=1,2$), \lean{ScalingLaw} (scaling for $m=1$ exact, $m=2$ with
$1/d_1^5$ error), \lean{Recurrence} (integer Chebyshev connection),
\lean{SubfactorBridge} (subfactor/Temperley--Lieb connection, 3 axioms),
\lean{GJSCircle} (GJS circle law), \lean{Verify} (numerical verification
via \lean{native\_decide}).

% ══════════════════════════════════════════════════════════════════
\section{Higher-order coefficients}
\label{sec:higher-order}

The main theorem (Theorem~\ref{thm:main}) uses only the first-order
coefficient~$c_1$ of the implicit function~$\psi$.  However, the
formalisation also develops the machinery needed to compute
\emph{all} Taylor coefficients $c_n$ of~$\psi(\delta) = \alpham +
\sum_{n \ge 1} c_n \delta^n$.  This section describes the algebraic
approach implemented in \lean{LagrangeCoefficients.lean}.

\subsection{The polynomial \texorpdfstring{$d_{\mathrm{poly}}$}{d\_poly} and root factoring}

Since $d(n, \lambda)$ is a polynomial in~$\lambda$ for each~$n$,
we introduce its algebraic counterpart:

\begin{definition}[\lean{LagrangeCoefficients.d\_poly}]
The polynomial $d_{\mathrm{poly}}(n) \in \R[X]$ is defined by the
same recurrence as~$d$:
\[
  d_{\mathrm{poly}}(0) = 1, \quad
  d_{\mathrm{poly}}(1) = X, \quad
  d_{\mathrm{poly}}(n{+}2) = X \cdot d_{\mathrm{poly}}(n{+}1) - X \cdot d_{\mathrm{poly}}(n),
\]
with the evaluation bridge
$d_{\mathrm{poly}}(n).\mathrm{eval}\,\lambda = d(n, \lambda)$.
\end{definition}

Since $\alpham$ is a root of $d_{\mathrm{poly}}(m{+}1)$
(Theorem~\ref{thm:vanish}), the factor theorem gives:

\begin{theorem}[\lean{d\_poly\_factor}]
\label{thm:d-poly-factor}
$d_{\mathrm{poly}}(m{+}1) = (X - \alpham) \cdot q_m$, where
$q_m = d_{\mathrm{poly}}(m{+}1) \mathbin{/\!\!/_{\mathrm{m}}} (X - \alpham)$
is the quotient after dividing by the monic linear factor.
\end{theorem}

The algebraic L'H\^opital rule then yields:

\begin{theorem}[\lean{q\_eval\_eq\_derivative}]
\label{thm:q-eval}
$q_m(\alpham) = d_{\mathrm{poly}}'(m{+}1)(\alpham) = d'(m{+}1,\alpham)$.
\end{theorem}

\begin{proof}
Differentiating $d_{\mathrm{poly}}(m{+}1) = (X - \alpham) \cdot q_m$
gives $d_{\mathrm{poly}}'(m{+}1) = q_m + (X - \alpham) \cdot q'_m$.
Evaluating at $\alpham$: $d_{\mathrm{poly}}'(m{+}1)(\alpham) = q_m(\alpham)$.
\end{proof}

This recovers the first-order coefficient purely algebraically:
$c_1 = d(m,\alpham)/q_m(\alpham)$, with no limit or L'H\^opital.

\subsection{Analyticity chain}

The higher-order coefficients require relating the Taylor
coefficients of~$\psi$ to those of~$d(k, \psi(\cdot))$.
Three Mathlib results provide the chain:

\begin{enumerate}[label=(\arabic*)]
  \item \emph{Polynomials are analytic.}
    Since $d(n, \cdot)$ is a polynomial, it is analytic at every
    point (\lean{AnalyticOnNhd.eval\_polynomial}).
  \item \emph{The ratio $g(\alpha) = d(m{+}1,\alpha)/d(m,\alpha)$ is analytic}
    near~$\alpham$, by \lean{AnalyticAt.div} (the denominator
    $d(m, \alpham) \ne 0$).
  \item \emph{The inverse function $\psi = g^{-1}$ is analytic} at
    $g(\alpham) = 0$, by the analytic inverse function theorem
    (\lean{analyticAt\_localInverse}), since $g'(\alpham) \ne 0$.
\end{enumerate}

\subsection{Fa\`a di Bruno via \texorpdfstring{$\texttt{OrderedFinpartition}$}{OrderedFinpartition}}
\label{sec:fdb}

Mathlib's \lean{Mathlib.Analysis.Calculus.IteratedDeriv.FaaDiBruno}
provides the one-dimensional Fa\`a di Bruno formula directly in
terms of iterated derivatives, summing over ordered finite
partitions.  The general identity is:

\begin{theorem}[\lean{iteratedDeriv\_comp\_eq\_sum\_orderedFinpartition}]
\label{thm:fdb-general}
Let $g, f$ be sufficiently smooth.  Then for every $n \ge 1$:
\[
  \frac{d^n}{dx^n}(g \circ f)(x)
  \;=\; \sum_{c \,:\, \mathrm{OrderedFinpartition}\,n}
  g^{(|c|)}(f(x)) \cdot \prod_{j=1}^{|c|} f^{(c_j)}(x),
\]
where the sum ranges over all ordered partitions of~$\{1,\ldots,n\}$
into non-empty blocks, $|c|$ is the number of blocks, and
$c_j$ is the size of the $j$-th block.
\end{theorem}

For small $n$, this specialises to explicit formulas already in Mathlib:

\begin{itemize}[nosep]
  \item $n = 2$: $\;(g \circ f)'' = g''(f) \cdot (f')^2 + g'(f) \cdot f''$
    \quad(\lean{iteratedDeriv\_comp\_two}),
  \item $n = 3$: $\;(g \circ f)''' = g'''(f) \cdot (f')^3 + 3\,g''(f) \cdot f' \cdot f'' + g'(f) \cdot f'''$
    \quad(\lean{iteratedDeriv\_comp\_three}).
\end{itemize}

\noindent
The analytic composition route via \lean{HasFPowerSeriesAt.comp}
and \lean{FormalMultilinearSeries.comp} remains available as an
alternative (it sums over \lean{Composition}\,$n$ rather than
\lean{OrderedFinpartition}\,$n$), but the iterated-derivative
formulation connects directly to the Leibniz recursion below.

\subsection{The Leibniz recursion}

The linearity of $F(\delta, \alpha) = d(m{+}1,\alpha) - \delta \cdot d(m,\alpha)$
in~$\delta$ yields a key simplification.  Differentiating
$d(m{+}1, \psi(\delta)) = \delta \cdot d(m, \psi(\delta))$
$n$~times and evaluating at $\delta = 0$:

\begin{theorem}[Leibniz recursion, \lean{leibniz\_recursion}]
\label{thm:leibniz}
For every $n \ge 1$,
\[
  \frac{d^n}{d\delta^n}\bigl[d(m{+}1, \psi(\delta))\bigr]\big|_0
  \;=\; n \cdot
  \frac{d^{n-1}}{d\delta^{n-1}}\bigl[d(m, \psi(\delta))\bigr]\big|_0.
\]
\end{theorem}

\begin{proof}
Apply the Leibniz product rule (\lean{iteratedDeriv\_mul}) to
$\delta \cdot d(m, \psi(\delta))$.  At $\delta = 0$, only the
$i = 1$ term survives in the binomial sum, contributing
$\binom{n}{1} \cdot 1 \cdot d^{(n-1)}[d(m,\psi)]|_0$.
\end{proof}

Combining the Leibniz recursion with the Fa\`a di Bruno formula
(Theorem~\ref{thm:fdb-general}) yields the complete coefficient
recursion:

\begin{theorem}[General recursion, \lean{fdb\_general\_recursion}]
\label{thm:general-recursion}
For every $m \ge 1$ and $n \ge 1$:
\[
  \sum_{c \,:\, \mathrm{OFP}(n)}
  d^{(|c|)}(m{+}1, \alpham) \cdot \prod_{j} \psi^{(c_j)}(0)
  \;=\; n \cdot
  \sum_{c \,:\, \mathrm{OFP}(n{-}1)}
  d^{(|c|)}(m, \alpham) \cdot \prod_{j} \psi^{(c_j)}(0),
\]
where $\mathrm{OFP}(k) = \mathrm{OrderedFinpartition}\,k$.
\end{theorem}

\begin{proof}
The left-hand side equals
$\frac{d^n}{d\delta^n}[d(m{+}1, \psi(\delta))]|_0$
by Theorem~\ref{thm:fdb-general}, and similarly for the right.
The identity then follows from Theorem~\ref{thm:leibniz}.
The Lean proof chains \lean{iteratedDeriv\_comp\_eq\_sum\_orderedFinpartition}
with \lean{leibniz\_recursion} and contains no \texttt{sorry}.
\end{proof}

This recursion isolates~$c_n$ from $c_1, \ldots, c_{n-1}$, as follows.

\begin{corollary}[Coefficient extraction, \lean{cn\_extraction}]
\label{cor:cn-extraction}
On the left-hand side of Theorem~\ref{thm:general-recursion}, the unique
ordered partition with a single block of size~$n$ (i.e.\ $|c| = 1$,
$c_1 = n$) contributes
$d'(m{+}1,\alpham) \cdot \psi^{(n)}(0) = d'(m{+}1,\alpham) \cdot n!\, c_n$.
Every other partition has $|c| \ge 2$, so every block size is ${<}\,n$
and the term depends only on $c_1, \ldots, c_{n-1}$.
Since $d'(m{+}1,\alpham) \ne 0$
(Theorem~\ref{thm:deriv-pos}), solving gives
\[
  \boxed{\;c_n \;=\; \frac{1}{n!\,d'(m{+}1,\alpham)}
  \Bigl[\,
    n \cdot R_{n-1}
    \;-\; L_n^{\ne}
  \Bigr]\;}
\]
where, writing $\psi^{(k)}(0) = k!\,c_k$:
\begin{align*}
  R_{n-1} &= \sum_{c \in \mathrm{OFP}(n{-}1)}
    d^{(|c|)}(m,\alpham)\,\prod_{j=1}^{|c|} (c_j)!\,c_{c_j}, \\
  L_n^{\ne} &= \!\!\sum_{\substack{c \in \mathrm{OFP}(n)\\ |c| \ge 2}}
    d^{(|c|)}(m{+}1,\alpham)\,\prod_{j=1}^{|c|} (c_j)!\,c_{c_j}.
\end{align*}
Both sums depend only on $c_1, \ldots, c_{n-1}$.
\end{corollary}

\begin{proof}[Lean proof sketch]
The proof proceeds in four steps:
\begin{enumerate}[label=(\roman*),nosep]
\item \emph{Sum splitting.}
  \lean{sum\_split} decomposes the OFP$(n)$ sum into
  partitions with $|c| = 1$ (the \lean{singleBlockSet}) and
  $|c| \ge 2$ (the \lean{multiBlockSet}).
\item \emph{Uniqueness.}
  \lean{eq\_singleBlock} proves the single-block partition is unique:
  any $c$ with $|c| = 1$ has part size~$n$ (from \lean{sum\_partSize})
  and a strictly monotone embedding $\mathrm{Fin}\,n \to \mathrm{Fin}\,n$,
  which must be the identity (\lean{strictMono\_fin\_eq\_id},
  via \lean{Fin.coe\_orderIso\_apply}).
\item \emph{Collapse.}
  \lean{singleBlockSet\_sum} evaluates the single-block sum to
  $d'(m{+}1,\alpham) \cdot \psi^{(n)}(0)$
  via \lean{singleBlock\_term}.
\item \emph{Rearrangement.}
  Since $d'(m{+}1,\alpham) \ne 0$, divide to isolate $\psi^{(n)}(0)$.
\end{enumerate}
The only remaining \texttt{sorry} in the entire file is the
final algebraic rearrangement (step~iv), which is a
\lean{field\_simp} / \lean{linarith} computation.
\end{proof}

The atomic ingredients at $\alpham = 4\cos^2\theta$,
$\theta = \pi/(m{+}2)$:

\medskip
\begin{center}
\begin{tabular}{@{}ll@{}}
\toprule
Quantity & Closed form \\
\midrule
$d(k,\alpham)$ & $\sin((k{+}1)\theta)/\sin\theta$ \\
$d(m,\alpham)$ & $(2\cos\theta)^m$ \\
$d'(m{+}1,\alpham)$ & $(2\cos\theta)^m \cdot (m{+}2)/(4\sin^2\theta)$ \\
$d^{(r)}(k,\alpham)$ & \lean{d\_deriv\_iter} via the Chebyshev connection \\
\bottomrule
\end{tabular}
\end{center}
\medskip

\noindent\emph{Verification at $n = 1$.}
$\mathrm{OFP}(1)$ has a single element (the single block)
and $\mathrm{OFP}(0)$ is the empty partition.
The formula gives
$c_1 = d(m,\alpham)/d'(m{+}1,\alpham)
     = 4\sin^2\theta/(m{+}2)$,
matching \lean{first\_order\_coeff}
(proved as \lean{taylor\_coeff\_one}).

\smallskip
\noindent\emph{Verification at $n = 2$.}
$\mathrm{OFP}(2)$ has two elements: the single block $\{1,2\}$ and
the pair of singletons $\{1\},\{2\}$.
So $L_2^{\ne} = d''(m{+}1,\alpham) \cdot c_1^2$ and
$R_1 = d'(m,\alpham) \cdot c_1$, giving
\[
  c_2 = \frac{d'(m,\alpham)\cdot c_1 - \tfrac{1}{2}\,d''(m{+}1,\alpham)\cdot c_1^2}
             {d'(m{+}1,\alpham)}.
\]
After substituting Chebyshev values this yields
Definition~\ref{def:c2}:
$c_2 = \sin^2\theta\,(2(m{+}1)\cos^2\theta - 1) / ((m{+}2)^2 \cos^2\theta)$.

\smallskip
Since each $d^{(r)}(k, \alpham)$ is a polynomial in $\cos\theta$
and $\sin\theta$ with rational coefficients (from the Chebyshev
connection formula), induction on~$n$ using
Corollary~\ref{cor:cn-extraction} shows that every
$c_n \in \mathbb{Q}(\cos\theta, \sin\theta)$.

\subsection{Second-order coefficient}

\begin{definition}[\lean{second\_order\_coeff}]
\label{def:c2}
\[
  c_2(m) \;=\;
  \frac{\sin^2\theta \cdot (2(m{+}1)\cos^2\theta - 1)}
       {\cos^2\theta \cdot (m{+}2)^2},
  \qquad \theta = \frac{\pi}{m+2}.
\]
\end{definition}

\begin{remark}
For $m = 1$: $c_2 = 0$ (the tridiagonal expansion terminates
at first order, consistent with the exact formula
$\psi(\delta) = 1 + \delta$).
For $m = 2$: $c_2 = 1/8$, giving
$\psi(\delta) = 2 + (\sqrt{2}/2)\,\delta + \delta^2/8 + O(\delta^3)$.
\end{remark}

The Lean formalisation verifies explicit values of~$c_2$ for small~$m$:

\medskip
\begin{center}
\begin{tabular}{@{}ccccl@{}}
\toprule
$m$ & $\theta$ & $\alpham$ & $c_2(m)$ & Lean status \\
\midrule
$1$ & $\pi/3$ & $1$ & $0$ & \lean{second\_order\_coeff\_m1} (\checkmark) \\
$2$ & $\pi/4$ & $2$ & $1/8$ & \lean{explicit\_m2\_c2} (\checkmark) \\
$3$ & $\pi/5$ & $(3{+}\sqrt{5})/2$ & $\sqrt{5}/25$ & \lean{explicit\_m3\_c2} (\texttt{sorry}) \\
$4$ & $\pi/6$ & $3$ & $13/216$ & \lean{explicit\_m4\_c2} (\checkmark) \\
\bottomrule
\end{tabular}
\end{center}
\medskip

\noindent
The $m = 3$ case involves algebraic manipulations with~$\sqrt{5}$
that \lean{norm\_num} does not close automatically; the remaining
\texttt{sorry} is purely computational.

\subsection{Palindromic polynomial reduction}
\label{sec:palindromic}

The Chebyshev identity
$2\cos\varphi \cdot \sin((m{+}2)\varphi)
 = \sin((m{+}3)\varphi) + \sin((m{+}1)\varphi)$
rewrites the implicit equation
$\delta = 2\cos\varphi \cdot \sin((m{+}2)\varphi)/\sin((m{+}1)\varphi)$
as
\begin{equation}\label{eq:palindromic-trig}
  (\delta - 1)\,\sin\bigl((m{+}1)\varphi\bigr)
  = \sin\bigl((m{+}3)\varphi\bigr).
\end{equation}
Setting $u = e^{2i\varphi}$ (so that
$\alpha = (u^{1/2} + u^{-1/2})^2 = u + 2 + u^{-1}$)
and multiplying by $u^{m+3}$, equation~\eqref{eq:palindromic-trig}
becomes the \emph{palindromic} polynomial
\begin{equation}\label{eq:palindromic-poly}
  P(u) \;=\; u^{m+3} - \sigma\, u^{m+2} + \sigma\, u - 1 \;=\; 0,
  \qquad \sigma = \delta - 1.
\end{equation}
Since $P(1) = 0$ identically, we factor $P = (u-1)\,Q(u)$ where
\[
  Q(u) = u^{m+2} + (2-\delta)\bigl(u^{m+1} + u^m + \cdots + u\bigr) + 1.
\]
$Q$~is palindromic: $Q(u) = u^{m+2}\,Q(1/u)$.  Setting $v = u + u^{-1}$
(so that $\alpha = v + 2$), the palindromic symmetry halves the degree:
$Q$~reduces to a polynomial of degree $\lfloor(m{+}2)/2\rfloor$ in~$v$.
When $m+2$ is odd, $u = -1$ is an additional root of~$Q$, giving a
further factor of $(u+1)$ and reducing the effective degree by one.

\medskip
\begin{center}
\begin{tabular}{@{}cccl@{}}
\toprule
$m$ & $Q$ degree & degree in $v$ & status \\
\midrule
$1$ & $3$ & $1$ (linear) & $\alpha(\delta) = 1 + \delta$ \\
$2$--$3$ & $4$--$5$ & $2$ (quadratic) & square root \\
$4$--$5$ & $6$--$7$ & $3$ (cubic) & Cardano \\
$6$--$7$ & $8$--$9$ & $4$ (quartic) & radicals \\
$8$--$9$ & $10$--$11$ & $5$ (quintic) & Abel--Ruffini \\
\bottomrule
\end{tabular}
\end{center}
\medskip

\noindent\textbf{Explicit closed forms.}
For $m = 2$, the reduced equation is $v^2 + (2{-}\delta)\,v - \delta = 0$,
giving
\[
  \boxed{\alpha(\delta) = \frac{\delta + 2 + \sqrt{\delta^2 + 4}}{2}}
  \qquad (m = 2).
\]
In particular, $c_1 = \tfrac12$, $c_2 = \tfrac18$, all odd $c_{2k+1} = 0$
for $k \ge 1$, and
$c_{2k} = \binom{1/2}{k} / 4^k$ (generalised binomial coefficients)
--- a single binomial sum rather than the Bell-number partition sum
from Corollary~\ref{cor:cn-extraction}.

For $m = 3$, after factoring $u = -1$, the reduced equation is
$v^2 + (1{-}\delta)\,v - 1 = 0$, giving
\[
  \boxed{\alpha(\delta) = \frac{\delta + 3 + \sqrt{\delta^2 - 2\delta + 5}}{2}}
  \qquad (m = 3).
\]
At $\delta = 0$ this recovers $\alpha_3 = (3 + \sqrt{5})/2$
(the golden ratio plus one).

\subsection{Galois theory: arithmetic vs.\ functional}
\label{sec:galois}

The palindromic symmetry $Q(u) = u^{m+2}\,Q(1/u)$ is a genuine
Galois-theoretic constraint: the splitting field of~$Q(u)$ over
$\mathbb{Q}(\delta)$ sits in a short exact sequence
\[
  1 \;\longrightarrow\; (\mathbb{Z}/2)^d
  \;\longrightarrow\; \mathrm{Gal}\bigl(Q(u)/\mathbb{Q}(\delta)\bigr)
  \;\longrightarrow\; \mathrm{Gal}\bigl(\text{$v$-poly}/\mathbb{Q}(\delta)\bigr)
  \;\longrightarrow\; 1,
\]
where $d = \lfloor(m{+}2)/2\rfloor$ and the $(\mathbb{Z}/2)^d$ factor
comes from solving $u^2 - vu + 1 = 0$ for each root of the
$v$-polynomial.
Since $(\mathbb{Z}/2)^d$ is solvable, the full Galois group is
solvable if and only if $\mathrm{Gal}(\text{$v$-poly}/\mathbb{Q}(\delta))$
is solvable.  As the reduced degree $d$ reaches~5 only at $m = 8$,
this is the correct Abel--Ruffini threshold.

\medskip
\noindent\textbf{Arithmetic vs.\ functional obstruction.}
The \emph{constants} $\alpha_m = 4\cos^2(\pi/(m{+}2))$ lie in the
maximal real subfield $\mathbb{Q}(\zeta + \zeta^{-1})$ of the cyclotomic
field $\mathbb{Q}(\zeta)$, $\zeta = e^{2\pi i/(m+2)}$.  This is an
\emph{abelian} (hence solvable) extension of~$\mathbb{Q}$, so the
threshold values are always expressible in radicals, for every~$m$.
Similarly, each~$c_n$ lies in the same abelian field
$\mathbb{Q}(\cos\theta, \sin\theta)$.

The Abel--Ruffini obstruction applies only to the \emph{functional}
dependence: whether $\psi(\delta)$ can be expressed \emph{globally} in
radicals of~$\delta$ (i.e., the monodromy group of the $v$-polynomial
over $\mathbb{Q}(\delta)$).  For $m \le 7$, the degree is~$\le 4$
and the answer is yes (with explicit closed forms for $m \le 3$ given
above).  For $m \ge 8$, the monodromy can be non-solvable, though the
specific Chebyshev structure may still yield a solvable Galois group.

\medskip
\noindent\textbf{Computability.}
Regardless of $m$, the Taylor coefficients~$c_n$ are always computable:
the Lagrange--Leibniz recursion
(Theorem~\ref{thm:general-recursion}) expresses each~$c_n$ as a
rational function of $\cos\theta$, $\sin\theta$, and
$c_1, \ldots, c_{n-1}$, while the palindromic closed forms
(\S\ref{sec:palindromic}) give independent verification for
$m \le 7$.

\begin{remark}[Rationality]
Each $c_n$ lies in the field $\mathbb{Q}(\cos\theta, \sin\theta)$
$= \mathbb{Q}(\zeta + \zeta^{-1})$,
since all iterated derivatives
$d^{(r)}(k, \alpham)$ are polynomials in $\cos\theta$ and
$\sin\theta$ with rational coefficients (via the Chebyshev
connection formula), and the recursion involves only
rational operations.  This is an abelian extension of~$\mathbb{Q}$
of degree~$\varphi(m{+}2)/2$.
\end{remark}

\subsection{Linear recurrence for the Taylor coefficients}
\label{sec:recurrence}

Since $\psi(\delta)$ is algebraic---it satisfies the polynomial equation
$R(\psi(\delta), \delta) = 0$ obtained from the palindromic reduction---the
generating function $\sum_{n \ge 1} c_n \delta^n$ is an algebraic power
series.  By the theory of D-finite functions, its coefficients satisfy a
\emph{linear recurrence with polynomial coefficients} in~$n$:
\begin{equation}\label{eq:recurrence}
  \sum_{j=0}^{D} p_j(n) \, c_{n-j} = 0 \qquad (n \gg 0),
\end{equation}
where $D = \lfloor (m+2)/2 \rfloor$ is the algebraic degree from
\S\ref{sec:palindromic} and $p_j \in \mathbb{Q}[n]$ depend only on~$m$.

\smallskip
\noindent\textbf{Examples.}
For $m = 2$, the recurrence decouples into even and odd subsequences:
$(n+1)\,c_{n+1} = -\tfrac{n-1}{4}\,c_{n-1}$, recovering
$c_{2k} = \binom{1/2}{k}/4^k$ with all odd coefficients vanishing.
For $m = 3$: $(n+1)\,c_{n+1} = \bigl((n-1)\,c_{n-1} - 2n\,c_n\bigr)/4$.

\smallskip
\noindent\textbf{Comparison with the Lagrange recursion.}
The Lagrange--Leibniz recursion (Theorem~\ref{thm:general-recursion}) is
\emph{universal}---the same formula covers all~$m$---but involves
multinomial convolutions requiring $O(n^2)$ operations for $c_1, \ldots, c_n$.
The linear recurrence~\eqref{eq:recurrence} is \emph{specific} to each~$m$
but computes the same sequence in $O(n)$ operations after a one-time
$O(m^2)$ setup.  For the Lean formalisation, the universal recursion is
preferred (a single theorem covering all~$m$), while for numerical
computation of many coefficients, the linear recurrence is more efficient.

\begin{remark}[P-recursive, not C-recursive]
\label{rem:not-c-recursive}
The recurrence~\eqref{eq:recurrence} has \emph{polynomial} coefficients
$p_j(n)$, making $\{c_n\}$ a P-recursive (holonomic) sequence.  It is
\emph{not} C-recursive: the sequence $\{n\,c_n\}$ does not satisfy any
constant-coefficient linear recurrence of finite order.  In particular,
$c_n$ \textbf{cannot} be expressed as a finite exponential sum
\begin{equation}\label{eq:exp-sum-false}
  c_n \neq \frac{1}{n}\sum_{k=1}^r A_k\,\lambda_k^n
\end{equation}
for any constants $A_k, \lambda_k$, nor more generally as $(1/n)$ times
a sum of terms $P_k(n)\,\lambda_k^n$ with polynomial prefactors.
This was verified numerically: for both $m=2$ and $m=3$, the
sequence $\{n\,c_n\}$ fails every constant-coefficient recurrence test
at orders $2$ through~$6$.

The obstruction is structural.  The minimal polynomial
$R(\psi, \delta) = 0$ (degree~$\le 2$ in~$\psi$ for $m \le 3$) produces
a \emph{nonlinear} (quadratic) recurrence via the Cauchy product
$\sum_{j} c_j\,c_{n-j}$.  This can be linearised to the
P-recurrence~\eqref{eq:recurrence} by differentiating~$R$ and
eliminating~$\psi$, but the polynomial dependence on~$n$ is intrinsic.
Asymptotically, the algebraic branch-point singularity of~$\psi(\delta)$
gives $c_n \sim C\,n^{-3/2}\,\rho^{-n}$ (with $\rho$ the distance to
the nearest branch point), whose $n^{-3/2}$ factor is incompatible
with the $n^{-1}$ that any exponential sum~\eqref{eq:exp-sum-false}
would produce.
\end{remark}

\begin{remark}[The implicit function $g(\alpha)$]
\label{rem:g-function}
The function $g(\alpha) = \delta$ obtained by inverting the Chebyshev
parametrisation ($\alpha = 4\cos^2\!\phi$, $\delta = \sqrt{\alpha}
\cdot U_{m+1}(\sqrt{\alpha}/2) / U_m(\sqrt{\alpha}/2)$) is a rational
function of~$\alpha$.  It is \textbf{not} the ratio $d_{m+1}(\alpha)
/ d_m(\alpha)$ of consecutive determinants: the tridiagonal polynomial
$d_n(x)$ is defined with $x = 2\cos\phi$, so $d_n$ and $g$ live in
different variables ($x$ vs.\ $\alpha = x^2$).  For example:
$g(\alpha) = \alpha(\alpha{-}2)/(\alpha{-}1)$ for $m=2$,
and $g(\alpha) = (\alpha^2{-}3\alpha{+}1)/(\alpha{-}2)$ for $m=3$.
The Lagrange--B\"urmann formula inverts $g$ via
$h(\alpha) = (\alpha{-}\alpha_m)/g(\alpha)$ (a M\"obius transformation
for small~$m$), giving
\[
  c_n = \frac{1}{n}\bigl[\varepsilon^{n-1}\bigr]\,
  h(\alpha_m + \varepsilon)^n.
\]
Since $h$ is rational with a single pole (for $m = 2, 3$), the $n$-th
power $h^n$ has a pole of order~$n$, and the coefficient extraction
yields a finite sum of \emph{binomial-coefficient products}---not
simple exponentials.  For $m=2$:
\[
  c_n = \frac{1}{n}\sum_{k=0}^{n}\binom{n}{k}
  \frac{(-1)^{k+n-1}}{2^{k+n-1}}\binom{k+n-2}{n-1}.
\]
\end{remark}

% ══════════════════════════════════════════════════════════════════

\begin{thebibliography}{9}

\bibitem{notes}
K.~McShane and C.~Lancien,
\emph{PPT thresholds for random bipartite quantum states},
Notes, Institut Fourier, 2026.

\bibitem{jones}
V.~F.~R.~Jones,
Index for subfactors,
\emph{Invent.\ Math.}\ \textbf{72} (1983), 1--25.

\bibitem{gjs}
A.~Guionnet, V.~F.~R.~Jones, and D.~Shlyakhtenko,
Random matrices, free probability, planar algebras and subfactors,
\emph{Quanta of Maths}, Clay Math.\ Proc.\ \textbf{11} (2010), 201--239.

\bibitem{aubrun}
G.~Aubrun,
Partial transposition of random states and non-centered semicircular
distributions,
\emph{Random Matrices Theory Appl.}\ \textbf{1}(2) (2012), 1250001.

\end{thebibliography}

\end{document}
